% =============================================================================
% cap2_hardware/01_caracterizacao_pontos.tex
% Seção 2.1 -- Pontos de coleta no Campus 2 (v2.2 — segunda compactação)
% =============================================================================
\section{Pontos de coleta no Campus 2}
\label{sec:caracterizacao-pontos}

A discussão de hardware parte dos dez pontos de alimentação efetivamente operados no Campus 2 da USP São Carlos, não de um cenário idealizado. O levantamento foi conduzido em campo por inspeções cooperativas junto aos voluntários da atividade de extensão AEX Gatosdoc2, combinando ficha técnica, registro fotográfico georreferenciado e observação direta da rotina de alimentação --- procedimento alinhado às práticas de \emph{camera trapping} cooperativo, em que o conhecimento operacional dos cuidadores é informação primária~\cite{velez2022choosing,cordier2022africa}. A matriz consolidada P01--P10 (Apêndice~\ref{ap:matriz-pontos}) reúne os atributos de campo; aqui são sintetizados apenas aqueles que condicionam o problema de software do Capítulo~\ref{cap:pipeline}.

Três eixos delimitam esse problema. O primeiro é \emph{físico}: os abrigos são estruturas adaptadas pelos voluntários --- as ``casas de reciclagem'' construídas com caixas plásticas, telhas de fibrocimento, madeira recuperada e potes domésticos ---, dispostas no solo com geometria, altura e orientação variáveis. Trata-se, portanto, de infraestrutura não padronizada sobre a qual o módulo de captura deve ser retrofitável, sem obra civil e com baixo impacto visual. O segundo é \emph{energético e comunicacional}: sete dos dez pontos não dispõem de qualquer tomada de corrente alternada em raio operacionalmente útil ($\approx 20$\,m), o que torna inviável alimentação por rede elétrica como solução abrangente; em contrapartida, nove pontos recebem mais de quatro horas de incidência solar direta por dia (a exceção é o ponto P05, em corredor coberto). Wi-Fi USP e 4G alcançam nominalmente todos os pontos, mas a literatura adverte que disponibilidade nominal não equivale a estabilidade sustentada em ambiente externo arborizado~\cite{bruce2025largescale,velez2022choosing}, ressalva que motivará a opção por arquitetura \emph{offline-first}.

O terceiro eixo é \emph{ecológico} e tem consequência algorítmica direta. Por estarem dispostos no solo, os potes atraem rotineiramente fauna não-alvo recorrente: gambás (\estr{Didelphis albiventris}), quatis (\estr{Nasua nasua}), lagartos teiú, pombos, aves nativas e, criticamente, felinos silvestres como o gato-mourisco (\estr{Herpailurus yagouaroundi}), cuja presença em fragmentos urbanos da região central do estado de São Paulo é documentada em revisões recentes de fauna sinantrópica~\cite{goncalves2025wildcat,cove2022capital}. Essa biodiversidade altera o problema de software: não basta detectar presença diante do pote --- é preciso \textbf{discriminar a espécie} antes de submeter o quadro às etapas seguintes. A justificativa para o emprego de modelos generalistas como o \modeloSN~\cite{google_speciesnet_blog} e a pilha \modeloMD\,+\,PyTorch-Wildlife~\cite{hernandez2024pytorchwildlife,microsoft_cameratraps_repo} deriva diretamente dessa observação: filtrar o ruído biológico antes da reidentificação é condição necessária para que a re-ID de \estr{Felis catus} opere com taxas de falso positivo aceitáveis.

A leitura conjunta dos três eixos delimita o problema tratado no restante do capítulo: dimensionar uma infraestrutura de aquisição autônoma em energia na maioria dos pontos, tolerante a falhas de conectividade, capaz de operar em condições crepusculares e noturnas, discreta o suficiente para integrar-se às casinhas existentes e retrofitável sem alterar a infraestrutura física. Essas restrições definem os desafios que o pipeline do Capítulo~\ref{cap:pipeline} absorverá como parâmetros de entrada: variabilidade geométrica, iluminação heterogênea, intermitência de captura, presença sistemática de classes não-alvo e impossibilidade de pressupor canal de transmissão confiável.

Portanto, o hardware é tratado aqui como o contexto delimitador do problema; a análise dessas restrições físicas fornece os requisitos de entrada rigorosos para o pipeline de visão computacional, que constitui a verdadeira entrega de engenharia deste trabalho.